\begin{claim}{D.2.55}{Proved}{The incoming fifth differential does not change the bottom sigma group}
In the stable Atiyah--Hirzebruch spectral sequence of D.2.54,
\[
 d_5:E^5_{12,3}\longrightarrow E^5_{7,7}=\ZZ/4
       \quad\hbox{is zero},\qquad E^6_{7,7}=\ZZ/4.
 \tag{D.2.55a}
\]
More precisely, the bottom stable subgroup $\ZZ/4$ in the
eleven-stage injects after all twelve-cells are attached.
The doubled free $N'$-lift $p$ of D.2.22 satisfies the actual
identity
\[
                         (2p)\nu_{11}=0
                         \quad\hbox{in }\pi_{14}P_{11}.
                         \tag{D.2.55b}
\]
Here $P_{11}$ is the actual twice-shifted finite model of
D.2.22, before the twelve-cell relations are imposed.
Consequently only incoming $d_6$ and $d_7$ can change the
target in D.2.54 before its final $d_8$.
Neither those two images nor $d_8(H_0)$ is computed here;
no value of $h_Ml_4$ is asserted.
\end{claim}
\begin{proof}
\textbf{The finite stage and its complete coordinates.}
The evaluation $\Sigma^2E\to L=Y\langle7\rangle$ is
twelve-connected: its factorization through
$\Sigma\Omega L$ has successive connectivities twelve
and thirteen, as in the double-counit argument of D.2.43.
Thus it has a relative CW model with cells only in
dimensions at least thirteen.  The twice-shifted stable
model of D.2.22 through dimension twelve realizes
uniquely in the ordinary stable range and provides the
required stages $P_j$.  A later thirteen-cell can kill
the suspended $u_5^2$ sphere; it cannot alter the
twelve-stage calculation below.

All homotopy groups in the proof are stable unless an
ordinary cone calculation is specified.  The groups
$\pi_{11}P_{11}$ are also ordinary groups, since $P_{11}$
is six-connected.  Put
\[
 \begin{split}
 K&=(M_4\vee S^8_r)\cup_{i_7\eta_7+2r}e^9,
       \qquad M_4=M(\ZZ/4,7),\\
 K/S^7&=S^8_b\vee M(\ZZ/2,8)_r.
 \end{split}
 \tag{D.2.55c}
\]
This is the shifted core of D.2.22t2; the separate
$S^9_R$ has zero stable $\pi_{14}$.
The exact shifted component calculation gives
\[
 \pi_{11}K=\ZZ/8\{s_0\}\oplus\ZZ/2\{z\},\qquad
 2s_0=m_4,\quad z=r\nu_8,\quad
 \operatorname{pinch}(m_4)=2\nu_8.
 \tag{D.2.55d}
\]
After the ten- and eleven-cells, the complete old torsion
subgroup and free Hurewicz lattice are
\[
 \begin{split}
 V&=(\ZZ/2)^5\{s,z,\tau,w\eta_{10},j\eta_{10}\},\\
 h(p)&=4N',\qquad h(q)=2Q,\qquad
 h(x)=K_1,\qquad h(y)=K_5.
 \end{split}
 \tag{D.2.55e}
\]
These are D.2.22t9--t12, using the actual $a=b=0$
proved in D.2.31.  In particular $K_1,K_5$ attach
trivially and can be taken as free eleven-spheres.
Changing $x,y$ with the same full Hurewicz coordinates
adds elements of $V$ and disappears on doubling.

The bottom subgroup is
\[
 \operatorname{im}(\pi_{14}^sS^7\to\pi_{14}^sK)
                     =\ZZ/4\{i_7\sigma\}.
                     \tag{D.2.55f}
\]
The Moore attachment quotients the seventh stable stem
$\ZZ/16$ by four.  The additional nine-cell does not
kill anything further in this group: its possible
$\pi_{14}$ composite is
$(i_7\eta_7+2r)\nu_8^2=0$, since
$\eta\nu^2=0$ and $2\nu^2=0$.
The map $\pi_{14}^sK\to\pi_{14}^sP_{11}$ is injective:
ten- and eleven-cell attaching spheres have zero
$\pi_{14}^s$, by $\pi_5^s\mathbb S=\pi_4^s\mathbb S=0$.
The corresponding stable cofiber exact sequences justify
this old-subgroup injection without assumptions about
the later quotient $P_{11}/S^7$.

\textbf{The doubled $N'$-lift.}
Orient the $N'$ eleven-cell so its attaching map is
$\alpha\nu_7$, where $\alpha:S^7\to P_{11}$ is the
bottom Moore map.  It supplies an actual map
\[
             j_\nu:C\nu_7=S^7\cup_{\nu_7}e^{11}
                         \longrightarrow P_{11}.
                         \tag{D.2.55g}
\]
There is a unique $v_8\in\pi_{11}C\nu_7$ of top
degree eight: the pinch image is $8\ZZ$ and
$\pi_{11}S^7=0$.  Both $2p$ and $j_\nu v_8$ have
integral Hurewicz $8N'$.  Their difference lies in $V$.
Its second- and fourth-square cone detectors vanish.
For $2p$ this follows from the degree-two cone map;
for $j_\nu v_8$ it follows by naturality, since
$C\nu_7$ has no cohomology in degrees eight or ten.
The five detector vectors of D.2.22t10 are injective
on $V$.  Therefore
\[
                            2p=j_\nu v_8.
                            \tag{D.2.55h}
\]
This fixes the coextension, rather than discarding an
unmeasured old-cell difference.

The opened ordinary sphere convention gives
\[
       \{\nu_7,8\iota_{10},\nu_{10}\}=\{4\sigma'\}.
       \tag{D.2.55i}
\]
Indeed the author's program in \cite{FGHStricklandToda}
declares $\sigma'''=\{\nu_5,8\iota_8,\nu_8\}$ and
$E\sigma'''=2\sigma''$, $E\sigma''=2\sigma'$,
attributing them to Toda Lemmas 5.13--5.14.
Twice suspending the same ordinary cone choices gives
\textup{(D.2.55i)}; its indeterminacy
$\nu_7\pi_{14}S^{10}+\pi_{11}S^7\nu_{11}$ is zero.
Thus the actual cone composition is
\[
 (2p)\nu_{11}=j_\nu(v_8\nu_{11})
             =\alpha(4\sigma')=0.
             \tag{D.2.55j}
\]
For the last equality, $\alpha$ factors through $M_4$,
and $(4\iota_7)\sigma'=4\sigma'$ because $S^7$ is an
$H$-space.  Its Whitehead-square Hilton term is zero.
This proves the actual assertion \textup{(D.2.55b)}
before any twelve-cell imposing $2p=0$ is attached.

We also need the pure Moore composite
\[
                            m_4\nu_{11}=0.
                            \tag{D.2.55k}
\]
The cone formula places it in the bottom image of
$\{4\iota_7,2\nu_7,\nu_{10}\}$.
The defined bracket $\{4\iota_3,\nu',\nu_6\}$ lies
in $\pi_{10}(S^3)_{(2)}=0$; four suspensions put zero
in the former bracket.  Its only indeterminacy is
$\langle4\sigma'\rangle$, whose bottom image in $M_4$
is zero by the same $S^7$ argument.
The exact sphere vanishings and $E^4\nu'=2\nu_7$
are the opened formulas in \cite{FGHStricklandToda},
also used in D.2.53.

\textbf{All remaining orders, with the even normalization retained.}
Equations \textup{(D.2.55d), (D.2.55k)} imply that
$s_0\nu_{11}$ and $z\nu_{11}$ have order at most two.
Their images under the original quotient in
\textup{(D.2.55c)} have coordinates
\[
  (\nu_8^2,*),\qquad(0,i_8\nu_8^2),
  \qquad
  \pi_{14}^sM(\ZZ/2,8)=\ZZ/2\{i_8\nu_8^2\}.
  \tag{D.2.55l}
\]
The unspecified coordinate is immaterial: these two
vectors are independent.  Their span therefore misses
the bottom subgroup \textup{(D.2.55f)}.  This conclusion
passes to $P_{11}$ by the old-$K$ injection already
proved, without asserting that $P_{11}/S^7$ retains
the original quotient's group.
Also
\[
 \tau\nu_{11}=w\eta_{10}\nu_{11}
              =j\eta_{10}\nu_{11}=0,
 \tag{D.2.55m}
\]
since $\pi_{14}^sM(\ZZ/2,9)=\pi_5^s(\mathbb S/2)=0$
and $\eta\nu=0$.

For $q$ we must retain the even coordinate change
allowed in D.2.22t11.  Write the original $Q$ attachment
as $R\eta_9+2c n$, where $c\in\{0,1\}$ and
$n=\alpha\nu_7$.  Sliding $Q$ by $2cN'$ gives
$Q'=Q-2cN'$ attached by $R\eta_9$.  The resulting
actual $C\eta_9$ map provides $q_0$ with
\[
                  h(q_0)=2Q-4cN',
       \qquad q-q_0-cp\in V.
       \tag{D.2.55n}
\]
Pinch gives an isomorphism
$\pi_{14}C\eta_9\to\pi_{14}S^{11}=\ZZ/8$:
the surrounding fourth and fifth stable stems vanish.
Thus $q_0\nu_{11}$ has order four, since its pinch is
$2\nu_{11}$.  Equations \textup{(D.2.55e), (D.2.55j)}
now imply
\[
                            4q\nu_{11}=0.
                            \tag{D.2.55o}
\]
This accounts for the possible even $n$ correction.
The original $Q$ pinch of $q\nu_{11}$ is still
$2\nu_{11}$ and its other eleven-cell coordinates
are zero.  Likewise $2x\nu_{11},2y\nu_{11}$ have
order four and independent coordinates $2\nu_{11}$
on the free $K_1,K_5$ eleven-spheres.

\textbf{Every twelve-cell and the bottom intersection.}
The entire shifted attachment list D.2.22t13--t15,
with $a=b=0$, now has $\nu_{11}$-images generated by
\[
 s\nu_{11},\quad s\nu_{11}+z\nu_{11},\quad
 q\nu_{11}+s\nu_{11},\quad
 2x\nu_{11},\quad2y\nu_{11}.
 \tag{D.2.55p}
\]
The respective nonrepeated cells are $B_5,T_{10},B_4,
B_3,B_7$ in the original dimension-ten notation.
The $B_6,Z_2,B_2$ composites are zero; the other cells
are zero or repeat this list.  In particular all free
integral directions have been retained.

Suppose a linear combination of
\textup{(D.2.55p)} is in the bottom subgroup
\textup{(D.2.55f)}.  Collapse the lower skeleton and
project onto the $Q,K_1,K_5$ eleven-spheres.
The last three coefficients must each be divisible
by four, since their respective coordinates are
$2\nu\in\ZZ/8$.
Those whole terms vanish by
\textup{(D.2.55o)} and the order bounds above.
What remains is a combination of $s\nu$ and $z\nu$.
It comes from $K$, which injects on old $\pi_{14}^s$,
and its two quotient coordinates in
\textup{(D.2.55l)} force both coefficients even.
It too is zero.  Therefore
\[
 \langle\hbox{all twelve-cell }\nu\hbox{-composites}\rangle
   \quad\cap\quad
 \operatorname{im}(\pi_{14}^sS^7\to\pi_{14}^sP_{11})=0.
 \tag{D.2.55q}
\]

The stable cofiber sequence for the simultaneous
twelve-cell attachment identifies its kernel on
$\pi_{14}^s$ with exactly the subgroup on the left:
$\pi_{14}^s(\bigvee S^{11})$ is the direct sum of
the groups $\ZZ/8\{\nu_{11}\}$.
Thus the complete bottom $\ZZ/4$ injects into
$\pi_{14}^sP_{12}$.
In the cellular exact couple an incoming $d_5$ at
$(7,7)$ is a bottom stable class killed at this
twelve-stage after its lower-filtration components
have been removed
\cite[Part III, Section 7]{EH5Adams1974AHSS}.
The stronger injection just proved excludes every
such class.  This proves \textup{(D.2.55a)} without
identifying the page-five source with raw homology.
Thirteen-cell cellular boundaries may further
restrict that source and do not change the conclusion.
D.2.54 then leaves only incoming $d_6,d_7$, as stated.
\end{proof}
